\documentclass{article}
\usepackage{amsmath,hyperref}
\begin{document}
\section*{A received supplier byte is not a reservation of its route}
\paragraph{Lift and comparison.}
PR 858 made early positive callbacks and byte accounting native but supplied
the unavailable supplier as a test flag. This follow-up joins that exact
warehouse archive operation to the selected SubT visibility, range model,
breadcrumb-registration, dispatch, and receive callbacks used for PR 738.
The five vertices are projected from shortest-path weights on the published
Cave Circuit 01 graph. Pose updates, a common serial clock, and the output
archive's framing and authorization are fixture interfaces. Range is
deliberately 72\,m rather than the stock 500\,m. The test does not establish
a physical radio frequency law, field geometry, or an independent mission.
The selected dispatch code still charges an attempt before it can reject
delivery. The preparation and independent warehouse motions retain their
declared fixture semantics. Only P/Q sends to the old duty's receiver are
routed through this native visibility bridge.

\paragraph{Same action interface, different lawful future.}
Before any static registration, P reaches the old receiver and Q does not.
A moving-name announcement does not change that result. A near static relay
adds a cheaper route through a 75\,m hop for P (which fails the 72\,m cutoff)
but gives Q a viable route. A far static relay keeps P viable and opens Q.
The old header arrives before any test registration. In an ordered
schedule, the near event happens before the two early byte attempts;
Q then has the actual positive callback and finishes the pending return.
Both suppliers can be considered if both produce callbacks; the selector
chooses only a licensed and affordable one.

The key pair keeps the application prefix, authorization, copied result,
early callback from P, metered debit, and preparation reply fixed. In one
history P remains usable; in the other, the near static registration is
processed after the early callbacks and before the suffix. The suffix
attempt in the latter is charged in full and does not arrive. Making the
future registration possible enlarges the cell of physical continuations;
no supplier with a bound early byte then survives the intersection of its
menus. The same admitted-row checker withholds the suffix in both histories
under that widened scope, preserving independent work but not the old
deadline. A late far registration, by contrast, leaves P admissible.
This is a statement about the tested continuation families, not a
certificate that the stock broker can freeze or attest its topology.

\paragraph{Exact accounting and residual.}
With P's prior debit $d$ and header of 114 bytes, its early byte plus
full suffix needs $d+114+1+173\le320$; reusing the byte needs
$d+114+1+172\le320$. The useful boundary on the stable P route is
$d=33$. The comparison shares both the early observations and their
costs. A direct scheduler using the same future support and codec ties.
Predicate agreement before a new registration does not transport to the
changed continuation without either an enforced hold, later binding, or a
newly justified evidence path. Neither inserting a topology flag in the
checker nor relabeling a callback proves any such mechanism exists.
\end{document}
